\chapter{FGA representability of the Picard scheme}
\label{chap:Picard_FGAPicRepresentability}
% archon:covers AlgebraicJacobian/Picard/FGAPicRepresentability.lean

% STRATEGY NOTE (iter-174). This chapter is the A.2.c assembly sub-row of
% Route~A (positive-genus arm of nonempty_jacobianWitness). It wires the
% Quot-scheme engine (\cref{chap:Picard_QuotScheme}, A.2.b) and the
% \'etale-sheafified relative Picard functor (\cref{chap:Picard_RelPicFunctor},
% A.1.c) into the FGA representability statement for \(\Pic_{C/k}\), following
% Kleiman, ``The Picard scheme'' \S 4 (\cite{kleiman-picard}). It is a small
% \emph{assembly} chapter: every nontrivial input is a \(\) pointer to a
% sibling chapter, and the work here is to thread the inputs through
% Grothendieck's existence theorem (\cite[\S 4, Main Theorem
% th:main]{kleiman-picard}) and its specialisation to a complete variety over a
% field (\cite[\S 4, Cor.~cor:algsch]{kleiman-picard}). Identity-component
% material (\Picz) is A.3 and lives in a separate sibling chapter --- it is
% \emph{not} scoped into this chapter.

% NOTE (planner directive iter-174). This chapter must produce the
% \emph{representable Picard scheme} \(\Pic_{C/k}\) as a \(k\)-group scheme,
% locally of finite type, that represents the \'etale sheafification
% \(\Pic^\sharp_{(C/k)\et}\) defined in A.1.c. The three deliverables are
% (i) a line-bundle / Quot translation expressed via the Abel map
% \(\mathrm{Div}_{C/k} \to \Pic_{C/k}\) (Kleiman \S 3, \S 4) which factorises through the
% Hilbert scheme \(\Hilb_{C/k} = \Quot_{\mathcal{O}_C/C/k}\), hence through the
% Quot engine of A.2.b; (ii) the existence theorem itself, in the
% field-of-definition specialisation \(S = \Spec k\), \(X = C\) a smooth proper
% geometrically integral curve; (iii) the group-scheme refinement carried by
% the tensor product of invertible sheaves on \(C \times_k T\). The Lean
% encoding is delegated to a Lean file
% \texttt{AlgebraicJacobian/Picard/FGAPicRepresentability.lean} whose
% construction is described in the final section but not written in this
% chapter.

\section{Setup and motivation}
\label{sec:fga_pic_setup}

Throughout this chapter, fix a base field \(k\) and a smooth proper geometrically
integral curve \(C\) over \(k\) (the curve carrying the Jacobian). The relative
Picard presheaf \(\Pic^\sharp_{C/k} : (\Sch/k)^{op} \to \Set\) of
\cref{chap:Picard_LineBundlePullback} has already been packaged --- as a
set-valued functor at A.1.b, and refined to its abelian-group-valued
\'etale-sheafified version \(\Pic^\sharp_{(C/k)\et}\) at
\cref{chap:Picard_RelPicFunctor} (A.1.c). The Quot scheme machinery of
\cref{chap:Picard_QuotScheme} (A.2.b) supplies the representing engine: for
\(C/k\) projective and flat with integral geometric fibers, the Hilbert scheme
\(\Hilb_{C/k}\) and its open subscheme \(\mathrm{Div}_{C/k}\) of relative effective
divisors are both representable by quasi-projective \(k\)-schemes. Grothendieck's
existence theorem for \(\Pic\) assembles these two ingredients via the Abel
map \(\mathrm{Div}_{C/k} \to \Pic^\sharp_{C/k}\) to produce a representing scheme,
the \emph{Picard scheme} \(\Pic_{C/k}\). The present chapter packages this
assembly: a line-bundle / Quot correspondence, the FGA representability
theorem, and the abelian-group-scheme refinement.

\medskip

\textbf{The sheafification route (run-0008 rewire).} Kleiman \S 4 represents
the \'etale-sheafified functor \(\Pic_{(C/k)\et}\). This chapter instead
works throughout under the standing hypothesis that \emph{\(C\) has a
\(k\)-rational point} (\cref{def:has_rational_point}) and states
representability against the \emph{plain} relative Picard functor
\(\Pic^\sharp_{C/k}\) of \cref{def:rel_pic_sharp}. This is justified by
Kleiman \S 2, Thm.~2.5 (\texttt{th:comp}): if \(f : X \to S\) has a section
and \(\mathcal O_{S'} = f'_* \mathcal O_{X'}\) holds for all base changes
(automatic for our proper geometrically integral curve, by cohomology and
base change), then the comparison maps
\[
  \Pic_{X/S}(T) \longrightarrow \Pic_{(X/S)\et}(T) \longrightarrow
  \Pic_{(X/S)\fppf}(T)
\]
are bijective for every \(T\). Hence under the rational-point hypothesis the
plain relative functor \emph{is} the \'etale sheaf, and no \'etale
sheafification argument is needed here: this chapter's route bypasses it
entirely via the rational point. (Mathlib v4.31 does in fact provide an
\'etale Grothendieck topology on schemes,
\texttt{AlgebraicGeometry.etaleTopology}
[\texttt{Mathlib/AlgebraicGeometry/Sites/Etale.lean}]; the rational-point
route is a deliberate simplification, not a Mathlib-gap workaround.)
Without a section the plain functor is not even a Zariski sheaf in
general, so every representability assertion of this chapter carries the
rational-point hypothesis; the Jacobian consumers downstream
(\cref{chap:AbelJacobi}) work with a pointed curve \((C, P)\), which
supplies it.

% SOURCE: [Kleiman], ``The Picard scheme'', \S 4 (FGA Explained Ch.~9 \S 9.4);
% arXiv:math/0504020, p.~26. Read from
% references/kleiman-picard-src/kleiman-picard.tex.

\section{Line bundles and the Quot / Hilbert scheme}
\label{sec:fga_pic_line_bundle_quot}

The bridge from line bundles on \(C \times_k T\) to the Quot scheme is the
Abel map: a relative effective divisor \(D \subset C \times_k T\) over \(T\)
(an \(S\)-flat closed subscheme whose ideal is invertible) determines a closed
subscheme of \(C \times_k T\) flat over \(T\), hence a \(T\)-point of the Hilbert
functor \(\Hilb_{C/k}(T)\), which is the special case \(\Quot_{\mathcal{O}_C/C/k}\)
of the Quot functor of \cref{chap:Picard_QuotScheme}; and it determines an
invertible sheaf \(\mathcal{O}_{C \times_k T}(D)\) on the relative curve,
representing a class in \(\Pic^\sharp_{C/k}(T)\).

\begin{lemma}\leanok
[Line bundle / Quot correspondence via the Abel map]
  \label{lem:line_bundle_quot_correspondence}
  \lean{AlgebraicGeometry.Scheme.PicScheme.abelMapWitness,
    AlgebraicGeometry.Scheme.PicScheme.abelMap_app_mk}
  \uses{def:has_abel_map, def:inst_has_abel_map,
        def:div_functor_carrier, def:pic_sharp_carrier}
  % HONESTY NOTE (record of the wave-7 substantive discharge of Sorry C,
  % run 0011): this node's Lean content is the NATURAL-TRANSFORMATION
  % half of the Abel map only. It does NOT claim, and is not used to
  % derive, representability of the source Div_{C/k} by an open
  % subscheme of the Hilbert scheme (Kleiman th:repDiv), nor
  % representability of the Abel map itself by a morphism of k-schemes;
  % Div_{C/k} remains the Type-valued carrier presheaf of
  % def:div_functor_carrier throughout (see subsec:sorry_has_div_functor).
  % The scheme-representability content of the surrounding narrative
  % (Div is representable, hence so is the Abel map) is recorded as
  % context/history below the proof, kept separate and NOT covered by
  % \leanok.
  % SOURCE: [Kleiman], ``The Picard scheme'', \S 3, Def.~dfn:Abel (the Abel
  % map), and Thm.~th:repDiv (representability of relative effective
  % divisors by an open subscheme of \(\IHilb_{X/S}\)); arXiv:math/0504020,
  % pp.~25, 23 (read from
  % references/kleiman-picard-src/kleiman-picard.tex, L1837-L1841 + L2135-L2145).
  % SOURCE QUOTE:
  % "\begin{thm}\label{th:repDiv}
  %   Assume \(X/S\) is projective and flat.  Then \(\mathrm{Div}_{X/S}\) is
  %  representable by an open subscheme \(\IDiv_{X/S}\) of the Hilbert scheme
  %  \(\IHilb_{X/S}\).
  %   \end{thm}"
  % SOURCE QUOTE:
  % "\begin{dfn}\label{dfn:Abel}
  %   The {\it Abel map\/} is the natural map of functors
  %      \(\)A_{X/S}(T)\:\mathrm{Div}_{X/S}(T)\to \Pic_{X/S}(T)\(\)
  %   defined by sending a relative effective divisor \(D\) on \(X_T/T\) to the
  %  sheaf \(\mc O_{\!X_T}(D)\).  The target \(\Pic_{X/S}\) may be replaced by any
  %  of its associated sheaves."
  \textit{Source: [Kleiman], ``The Picard scheme'', \S 3,
  Thm.~3.7 + Def.~4.6 (cf.\ [Nitsure], ``Construction of Hilbert
  and Quot Schemes'', \S 1, where \(\Hilb_{X/S} = \Quot_{\mathcal{O}_X/X/S}\)).}
  Let \(C/k\) be a smooth proper geometrically integral curve. There is a
  natural transformation of presheaves on \((\Sch/k)^{op}\)
  \[
    A_{C/k} = \mathrm{abelMap} \;\colon\; \mathrm{Div}_{C/k}
      \;\longrightarrow\; \Pic^\sharp_{C/k}
  \]
  (\(\mathrm{Div}_{C/k}\) the relative-divisor functor of
  \cref{def:div_functor_carrier}, \(\Pic^\sharp_{C/k}\) the relative
  Picard functor of \cref{def:pic_sharp_carrier}), with the following
  \emph{defining property}: at a \(T\)-point represented by a family
  \(\langle \mathcal F, q\rangle\) (\cref{def:div_family}: a locally free
  rank-\(1\) quotient \(q \colon \mathcal F \twoheadrightarrow
  \mathcal L\) on \(C \times_k T\) with \(D := \ker q\) an invertible
  ideal sheaf, cutting out the relative effective divisor \(D \subseteq C
  \times_k T\)),
  \[
    A_{C/k}(T)\bigl(\langle \mathcal F, q\rangle\bigr)
      \;=\; -[\ker q] \;=\; [(\ker q)^{-1}]
      \;=\; \bigl[\mathcal{O}_{C \times_k T}(D)\bigr]
      \;\in\; \Pic^\sharp_{C/k}(T),
  \]
  the class of the associated invertible sheaf \(\mathcal{O}_{C \times_k
  T}(D)\), i.e.\ the additive inverse of the class of the ideal sheaf
  \(\ker q\) (Kleiman \S 3, Def.~dfn:Abel). This is the data carried by
  \(\mathrm{HasAbelMap}(C)\) (\cref{def:has_abel_map}), populated by the
  substantive witness \(\mathrm{abelMapWitness}(C)\) of
  \cref{def:inst_has_abel_map}.
\end{lemma}

\begin{proof}
  Write \(\mathrm{abelMapWitness}(C) = \mathrm{abelKernelNatTrans}(C)
  \mathbin{;} \mathrm{picNeg}(C)\), the composite of two natural
  transformations. The first, \(\mathrm{abelKernelNatTrans}(C) \colon
  \mathrm{Div}_{C/k} \to \Pic^\sharp_{C/k}\), sends
  \(\langle \mathcal F, q\rangle \mapsto [\ker q]\): it is well defined
  on the defining equivalence relation of \(\mathrm{Div}_{C/k}(T)\)
  (two families with proportional kernels have isomorphic kernels), and
  natural in \(T\) because kernel formation commutes with the flat base
  change \(g^* \colon C \times_k T \to C \times_k T'\) --- the comparison
  map \(g^*\ker q \to \ker(g^*q)\) is an isomorphism, since \(g^*q\) is
  again a surjection of locally free sheaves onto a locally free rank-1
  quotient. The second, \(\mathrm{picNeg}(C)\), is the (natural) additive
  inversion on the group-valued presheaf \(\Pic^\sharp_{C/k}\).
  Composing gives the natural transformation
  \(\mathrm{abelMapWitness}(C) \colon \langle \mathcal F, q\rangle
  \mapsto -[\ker q]\), and unfolding the two composites at a point gives
  the defining-property formula.
  \paragraph{Scope.} This proof establishes only the natural
  transformation \(A_{C/k}\) between the presheaves
  \(\mathrm{Div}_{C/k}\) and \(\Pic^\sharp_{C/k}\) of
  \cref{def:div_functor_carrier,def:pic_sharp_carrier}. It does
  \emph{not} establish, and does not use, representability of
  \(\mathrm{Div}_{C/k}\) by an open subscheme of the Hilbert scheme
  (Kleiman~\S 3, Thm.~th:repDiv, sketched in the remark below) or
  representability of \(A_{C/k}\) itself by a morphism of \(k\)-schemes;
  that scheme-level content is separate future work
  (\cref{subsec:sorry_has_div_functor}), consumed only by the
  (still-gated) Step~3 of \cref{thm:fga_pic_representability}.
\end{proof}

\begin{remark}
[Representability of \(\mathrm{Div}_{C/k}\) --- not yet formalised]
  \label{rem:div_functor_representability}
  % SOURCE QUOTE PROOF: (from references/kleiman-picard-src/kleiman-picard.tex,
  % L1843-L1866, proof of th:repDiv)
  % "Set \(H:=\IHilb_{X/S}\), and let \(W\subset X\x H\) be the universal
  % (closed) subscheme, and \(q\:W\to H\) the projection.  Let \(V\) denote the
  % set of points \(w\in W\) at which \(W\) is an effective divisor. Plainly \(V\)
  % is open in \(W\).  Set \(Z:=q(W-V)\).  Then \(Z\) is closed because \(q\) is
  % proper.  Set \(U:=H-Z\).  Then \(U\) is open, and \(q^{-1}U\) is an effective
  % divisor in \(X\x U\).  In fact, since \(q\) is flat, \(q^{-1}U\) is a relative
  % effective divisor in \(X\x U/U\).
  % [\dots]"
  Kleiman~\S 3, Thm.~th:repDiv, shows that \(\mathrm{Div}_{C/k}\) is
  representable by an open subscheme \(\mathrm{Div}_{C/k} \subseteq
  \Hilb_{C/k}\) of the Hilbert scheme (\(\Hilb_{C/k} =
  \Quot_{\mathcal{O}_C/C/k}\), \cref{chap:Picard_QuotScheme}): letting
  \(W \subseteq C \times_k \Hilb_{C/k}\) be the universal closed
  subscheme and \(q \colon W \to \Hilb_{C/k}\) the projection, the locus
  \(V \subseteq W\) where \(W\) is locally an effective Cartier divisor
  is open, so \(Z := q(W \setminus V)\) is closed (\(q\) proper) and
  \(\mathrm{Div}_{C/k} := \Hilb_{C/k} \setminus Z\) is the open
  subscheme representing relative effective divisors. This is
  \textbf{not} formalised: \(\mathrm{Div}_{C/k}\) remains the
  Type-valued carrier presheaf of \cref{def:div_functor_carrier}
  throughout the present chapter (\cref{subsec:sorry_has_div_functor}).
\end{remark}

\section{The FGA representability theorem}
\label{sec:fga_pic_representability}

We now state and outline the proof of Grothendieck's existence theorem
for the Picard scheme, in the specialisation needed for the Jacobian arm
of Route~A: \(S = \Spec k\), \(X = C\) a smooth proper geometrically integral
curve. The general theorem (Kleiman \S 4, Thm.~4.8) covers any
\(X/S\) projective Zariski-locally over \(S\), flat, with integral geometric
fibers; we record the general statement and then specialise.

\begin{theorem}\leanok
[FGA representability of the Picard scheme]
  \label{thm:fga_pic_representability}
  \lean{AlgebraicGeometry.Scheme.PicScheme.representable}
  \uses{lem:line_bundle_quot_correspondence, lem:smooth_proper_quotient,
        thm:quot_representable, def:rel_pic_sharp, def:has_rational_point,
        cor:flattening_stratification_curve, def:pic_sharp_representable,
        def:inst_pic_sharp_representable, lem:cech_flat_base_change}
  % SOURCE: [Kleiman], ``The Picard scheme'', \S 4, Main Thm.~th:main,
  % and Cor.~cor:algsch (specialisation to \(S = \Spec k\), \(X\) complete);
  % arXiv:math/0504020, pp.~26, 35 (read from
  % references/kleiman-picard-src/kleiman-picard.tex, L2155-L2166 + L2686-L2690).
  % SOURCE QUOTE:
  % "\begin{thm}[Main]\label{th:main}
  %   Assume \(f\:X\to S\) is projective Zariski locally over \(S\), and is flat
  %  with integral geometric fibers.
  %
  %  \tu{(1)} Then \(\IPic_{X/S}\) exists, is separated and locally of finite
  %  type over \(S\), and represents \(\Pic_{(X/S)\et}\).
  %
  %  \tu{(2)} If also \(S\) is Noetherian and \(X/S\) is projective, then
  %  \(\IPic_{X/S}\) is a disjoint union of open
  %   subschemes, each an increasing union of open quasi-projective
  %  \(S\)-schemes.
  %   \end{thm}"
  % SOURCE QUOTE:
  % "\begin{sbscor}\label{cor:algsch}
  %    Assume \(S\) is the spectrum of a field \(k\), and \(X\) is complete.  Then
  %  \(\IPic_{X/k}\) exists, represents \(\Pic_{(X/k)\fppf}\), and is a disjoint
  %  union of open quasi-projective subschemes.
  %   \end{sbscor}"
  \textit{Source: [Kleiman], ``The Picard scheme'', \S 4, Main Thm.~4.8
  + Cor.~4.18.3.}
  Let \(k\) be a field and \(C\) a smooth proper geometrically integral curve
  over \(k\) \emph{with a \(k\)-rational point}
  (\cref{def:has_rational_point}). Then the relative Picard functor
  \(\Pic^\sharp_{C/k}\) of \cref{def:rel_pic_sharp} --- which under the
  rational-point hypothesis coincides with its \'etale sheafification
  \(\Pic_{(C/k)\et}\), by Kleiman \S 2, Thm.~2.5 --- is representable by a
  \(k\)-scheme \(\Pic_{C/k}\), separated and locally of finite type over
  \(k\), which is a disjoint union of open quasi-projective
  \(k\)-subschemes, each an increasing union of open quasi-projective
  subschemes. The scheme \(\Pic_{C/k}\) is called the \emph{Picard scheme}
  of \(C/k\).
\end{theorem}

% SOURCE QUOTE PROOF: (from references/kleiman-picard-src/kleiman-picard.tex,
% L2168-L2366, proof of th:main; long and is broken into the four-step
% sketch below per the project's structural restatement convention. The
% verbatim quote of the opening reduction step is:)
% "By~\cite[(0, 4.5.5), p.~106]{EGAG}, it is a local matter on \(S\) to
% represent a Zariski sheaf on the category of \(S\)-schemes.  Moreover, it
% is also a local matter on \(S\) to prove that an \(S\)-scheme is separated
% and locally of finite type.  Hence, in order to prove (1), we may assume
% \(S\) is Noetherian and \(X/S\) is projective.
% [\dots]
% Plainly an \(S\)-scheme is separated if it is a disjoint union of
% separated open subschemes, or if it is an increasing union of
% separated open subschemes.  Hence (1) follows from (2)."
% (subsequent steps quoted inline in the project-notation proof below).
\begin{proof}
  
  We follow Kleiman~\S 4, Thm.~4.8, specialised to \(S = \Spec k\) and
  \(X = C\). The proof proceeds in four conceptual steps, preceded by a
  comparison step; the underlying algebraic engine is the Quot scheme of
  \cref{chap:Picard_QuotScheme} and the relative Picard functor of
  \cref{chap:Picard_RelPicFunctor}, together with the line-bundle / Quot
  correspondence of \cref{lem:line_bundle_quot_correspondence}.

  \emph{Step 0: reduce to the \'etale sheaf via the rational point.}
  By Kleiman \S 2, Thm.~2.5, since \(C \to \Spec k\) has a section
  (\cref{def:has_rational_point}) and \(\mathcal O_T = \pi_{T*}
  \mathcal O_{C \times_k T}\) holds for every \(k\)-scheme \(T\) (cohomology
  and base change for the proper geometrically integral fibers of
  \(C\times_k T/T\)), the comparison map
  \(\Pic^\sharp_{C/k}(T) \to \Pic_{(C/k)\et}(T)\) is bijective for all
  \(T\). It therefore suffices to represent the \'etale sheaf
  \(\Pic_{(C/k)\et}\), which is what Kleiman's theorem produces; the same
  scheme then represents \(\Pic^\sharp_{C/k}\).

  \emph{Step 1: stratify by Hilbert polynomial.}
  For each polynomial \(\phi \in \mathbb{Q}[n]\), let \(P^\phi \subseteq
  \Pic^\sharp_{(C/k)\et}\) be the \'etale subsheaf whose \(T\)-points are
  represented by invertible sheaves \(\mathcal{L}\) on \(C \times_k T\) with
  \(\chi(C_t, \mathcal{L}_t^{-1}(n)) = \phi(n)\) for all \(t \in T\) (the
  Hilbert polynomial relative to the chosen polarisation
  \(\mathcal{O}_C(1)\) on the projective curve \(C\)). The \(P^\phi\) form a
  disjoint open cover of \(\Pic^\sharp_{(C/k)\et}\) (openness via Kleiman's
  flat-base-change argument, \cref{lem:cech_flat_base_change} / Stacks Tag 02KH;
  cf.\ EGA III 7.9.11). It suffices to
  represent each \(P^\phi\) by an increasing union of open quasi-projective
  \(k\)-schemes.

  \emph{Step 2: bound twists by \(m\)-regularity.}
  Given \(\phi\), fix \(m\) large enough that the \(T\)-points of \(P^\phi\) are
  representable by \(\mathcal{L}\) on \(C \times_k T\) satisfying
  \(R^i \pi_{T*} \mathcal{L}(n) = 0\) for all \(i \ge 1\), \(n \ge m\). Twisting
  by \(\mathcal{O}_C(m)\) defines an automorphism of \(\Pic^\sharp_{(C/k)\et}\)
  that carries \(P^\phi\) to \(P^{\phi_0}_0\) for \(\phi_0(n) := \phi(m + n)\).
  Hence it suffices to represent each \(P^{\phi_0}_0\) by a quasi-projective
  \(k\)-scheme.

  \emph{Step 3: factor through the Abel map.}
  Consider the Abel map \(A_{C/k} : \mathrm{Div}_{C/k} \to \Pic^\sharp_{(C/k)\et}\)
  of \cref{lem:line_bundle_quot_correspondence}. The source \(\mathrm{Div}_{C/k}\)
  is representable by an open subscheme of \(\Hilb_{C/k}\) (and hence by an
  open subscheme of the Quot scheme of \cref{thm:quot_representable}), and
  is quasi-projective over \(k\). Pulling back along
  \(P^{\phi_0}_0 \hookrightarrow \Pic^\sharp_{(C/k)\et}\) yields an open
  subscheme \(Z := P^{\phi_0}_0 \times_{\Pic^\sharp_{(C/k)\et}} \mathrm{Div}_{C/k}\)
  of \(\mathrm{Div}_{C/k}\), which is again quasi-projective. The projection
  \(\alpha : Z \to P^{\phi_0}_0\) is a surjection of \'etale sheaves:
  given a \(T\)-point \(\lambda\) of \(P^{\phi_0}_0\) represented by an \'etale
  cover \(T' \to T\) and an invertible sheaf \(\mathcal{L}'\) on \(C \times_k
  T'\) with \(H^1(C_{t'}, \mathcal{L}'_{t'}) = 0\), the projective bundle
  \(\mathbb{P}(\mathcal{Q}) \to T'\) associated to
  \(\mathcal{Q} := \pi_{T'*}\mathcal{L}'\) is smooth over \(T'\); an \'etale
  cover \(T_1 \to T'\) admits a \(T'\)-map \(T_1 \to \mathbb{P}(\mathcal{Q})\)
  (EGA IV 17.16.3 (ii)) whose composition with the universal divisor map
  \(\mathbb{P}(\mathcal{Q}) \to Z\) produces the required \(T_1\)-point of \(Z\).
  Moreover, the first projection \(Z \times_{P^{\phi_0}_0} Z \to Z\) is
  smooth and proper.

  \emph{Step 4: descend by the smooth-proper quotient lemma.}
  Apply Kleiman~\S 4, Lem.~4.9: given a surjection \(\alpha : Z \to P\)
  of \'etale sheaves with \(Z\) quasi-projective, \(R := Z \times_P Z\)
  representable, and the first projection \(R \to Z\) smooth and proper,
  \(P\) is representable by a quasi-projective scheme. In our case
  \(P = P^{\phi_0}_0\), \(Z\) and \(R\) as above, so \(P^{\phi_0}_0\) is
  quasi-projective. Reassembling over \(\phi \in \mathbb{Q}[n]\) and
  over \(m\), the disjoint union of representable \(P^\phi\) pieces represents
  \(\Pic^\sharp_{(C/k)\et}\). The total scheme is separated (as a disjoint
  union of separated quasi-projectives) and locally of finite type. This
  is Kleiman~\S 4, Cor.~4.18.3 in the \(S = \Spec k\),
  \(X = C\) complete case.
\end{proof}

\subsection{The smooth-proper quotient lemma}
\label{subsec:fga_pic_qt}

The structural lemma underlying step 4 above is recorded separately
because it is consumed verbatim in the assembly proof: it is the
abstract content allowing the descent from \(Z = \mathrm{Div}^{\phi_0}_0\) to the
representing Picard piece.

\begin{lemma}\leanok
[Smooth-proper quotient of \'etale sheaves]
  \label{lem:smooth_proper_quotient}
  \lean{AlgebraicGeometry.Scheme.PicScheme.smoothProperQuotient}
  \uses{def:has_smooth_proper_quotient}
  % SOURCE: [Kleiman], ``The Picard scheme'', \S 4, Lem.~lm:qt;
  % arXiv:math/0504020, p.~31 (read from
  % references/kleiman-picard-src/kleiman-picard.tex, L2368-L2375).
  % SOURCE QUOTE:
  % "\begin{lem}\label{lm:qt}
  %   Let \(\alpha\:Z\to P\) be a map of \'etale sheaves, and set \(R:=Z\x_PZ\).
  %  Assume \(\alpha\) is a surjection, \(Z\) is representable by a
  %  quasi-projective \(S\)-scheme, \(R\) is representable by an \(S\)-scheme, and
  %  the first projection \(R\to Z\) is representable by a smooth and proper
  %  map.  Then \(P\) is representable by a quasi-projective \(S\)-scheme, and
  %  \(\alpha\) is representable by a smooth map."
  \textit{Source: [Kleiman], ``The Picard scheme'', \S 4, Lem.~4.9.}
  Let \(\alpha : Z \to P\) be a morphism of \'etale sheaves on \((\Sch/k)\).
  Set \(R := Z \times_P Z\). Suppose:
  \begin{enumerate}
    \item \(\alpha\) is a surjection of \'etale sheaves;
    \item \(Z\) is representable by a quasi-projective \(k\)-scheme;
    \item \(R\) is representable by a \(k\)-scheme;
    \item the first projection \(R \to Z\) is representable by a smooth
      and proper map.
  \end{enumerate}
  Then \(P\) is representable by a quasi-projective \(k\)-scheme, and \(\alpha\)
  is representable by a smooth map.
\end{lemma}

% NOTE (run-0008 truth audit): the LEAN statement of smoothProperQuotient is
% intentionally WEAKER than this blueprint statement: hypothesis (2)'s
% quasi-projectivity of Z is not expressible at the pinned Mathlib revision
% (no QuasiProjective vocabulary), and without it the naked statement is
% FALSE (a Hironaka-type free Z/2-action on a smooth proper non-projective
% 3-fold gives a smooth proper equivalence relation whose etale quotient
% sheaf is a non-scheme algebraic space). The Lean theorem therefore takes
% the conclusion-carrier class HasSmoothProperQuotient (def:has_smooth_proper_quotient)
% as a USE-SITE hypothesis, standing in for the missing quasi-projectivity;
% the former GLOBAL ⟨sorry⟩ instance (which asserted representability of
% EVERY presheaf target and was flatly false) was deleted in run 0008.
% Supplying the instance for the Abel-map slice (where Z is a genuinely
% quasi-projective open of Div_{C/k}) is part of the inst_has_pic_scheme
% obligation. Consequently the PROOF below is NOT formalised (no \leanok on
% the proof): it is the mathematical justification for the use-site
% instance, not for the Lean theorem body (which is a one-line extraction).

% SOURCE QUOTE PROOF: (from references/kleiman-picard-src/kleiman-picard.tex,
% L2377-L2417, proof of lm:qt; reproduced verbatim because it is the
% structural heart of step 4.)
% "To lighten the notation, let \(Z\) and \(R\) denote the
% corresponding schemes as well.  Since the structure map \(Z\to S\) is
% quasi-projective, it is separated; whence, the first projection
% \(Z\x_SZ\to Z\) is too.  But the first projection \(R\to Z\) is proper, and
% factors naturally through a map \(h\:R\to Z\x_SZ\).  Hence \(h\) is proper.
% But \(h\) is a monomorphism; that is, \(h\) is injective on \(T\)-points.
% Therefore, \(h\) is a closed embedding by \cite[8.11.5]{EGAIV3}.
%
% Plainly, for each \(S\)-scheme \(T\), the subset $R(T)\subset
% Z(T)\x_{S(T)}Z(T)$ is the graph of an equivalence relation.  Also, the
% map of schemes \(R\to Z\) is flat and proper, and \(Z\) is a
% quasi-projective \(S\)-scheme.  It follows that there exist a
% quasi-projective \(S\)-scheme \(Q\) and a faithfully flat and projective map
% \(Z\to Q\) such that \(R=Z\x_QZ\).  (In other words, a flat and proper
% equivalence relation on a quasi-projective scheme is effective.)
% [\dots]
% Since \(Z\to Q\) is flat, it is smooth if (and only if) its fibers are
% smooth.  But these fibers are, up to extension of the ground field, the
% same as those of \(R\to Z\).  And \(R\to Z\) is smooth by hypothesis.  Thus
% \(Z\to Q\) is smooth.
%
% It remains to see that \(Q\) represents \(P\).  First, observe that \(Z\to Q\)
% is (represents) a surjection of \'etale sheaves.
% [\dots]
% Since \(Z\to Q\) is a surjection of \'etale sheaves, \(Q\) is, in the
% category of \'etale sheaves, the coequalizer of the pair of maps
% \(R\rightrightarrows Z\) by Exercise~\ref{ex:epi} below, which is an
% elementary exercise in general Grothendieck topology.  By the same
% exercise, \(P\) too is the coequalizer of this pair of maps.  But, in any
% category, the coequalizer is unique up to unique isomorphism.  Thus \(Q\)
% represents \(P\)."
\begin{proof}
  By (ii) the structure map \(Z \to \Spec k\) is quasi-projective, hence
  separated, so the first projection \(Z \times_k Z \to Z\) is separated.
  By (iv) the first projection \(R \to Z\) is proper; it factors through
  \(h : R \to Z \times_k Z\) over the diagonal embedding (the second
  component is the second projection of \(R \to Z \times_P Z\)). The map
  \(h\) is therefore proper, and a monomorphism (because \(R\) is the graph
  of an equivalence relation on \(Z\), hence injective on \(T\)-points).
  By EGA IV 8.11.5, a proper monomorphism is a closed immersion: \(h\) is a
  closed embedding.

  Hence \(R \hookrightarrow Z \times_k Z\) is a flat (by (iv)) and proper
  equivalence relation on the quasi-projective \(k\)-scheme \(Z\). By the
  theory of flat-and-proper effective equivalence relations on
  quasi-projective schemes (Altman--Kleiman; the construction places
  the graph \(R \subseteq Z \times_k Z\) inside the universal subscheme of
  the Hilbert scheme \(\Hilb_{Z/k}\) and descends to a closed subscheme
  \(Q \subseteq \Hilb_{Z/k}\) representing the quotient), there exist a
  quasi-projective \(k\)-scheme \(Q\) and a faithfully flat, projective map
  \(Z \to Q\) with \(R = Z \times_Q Z\).

  The map \(Z \to Q\) is flat and proper; its fibers, up to field
  extension, coincide with those of \(R \to Z\), which are smooth by
  hypothesis. Hence \(Z \to Q\) is smooth.

  Finally, \(Z \to Q\) is a surjection of \'etale sheaves (an \'etale
  cover of any \(T\)-point of \(Q\) pulls back to a \(T'\)-point of \(Z\), by
  EGA IV 17.16.3 (ii) applied to \(Z \times_Q T \to T\)). In the category
  of \'etale sheaves, both \(Q\) and \(P\) are coequalisers of the pair
  \(R \rightrightarrows Z\); the coequaliser is unique up to unique
  isomorphism, so \(Q\) represents \(P\). Smoothness of \(\alpha : Z \to P\)
  is the smoothness of \(Z \to Q\) just shown.
\end{proof}

\section{Group-scheme structure}
\label{sec:fga_pic_group_structure}

The Picard scheme \(\Pic_{C/k}\) inherits an abelian group structure from
the tensor product of invertible sheaves on \(C \times_k T\). The relative
Picard presheaf \(T \mapsto \Pic(C \times_k T)/H_T\) is already an abelian
group at every \(T\) (the quotient of an abelian group by a subgroup),
functorially in \(T\); this lifts to the \'etale sheafification, and then
to a group-scheme structure on the representing scheme by Yoneda.

\begin{theorem}\leanok
[\(\Pic_{C/k}\) is a \(k\)-group scheme]
  \label{thm:pic_is_group_scheme}
  \lean{AlgebraicGeometry.Scheme.PicScheme.groupSchemeStructure}
  \uses{thm:fga_pic_representability, thm:relative_pic_quotient_well_defined,
        def:pic_scheme, def:rel_pic_sharp}

  % SOURCE: [Kleiman], ``The Picard scheme'', \S 2, Def.~df:Pfs (the relative
  % Picard functor takes values in abelian groups) and \S 4 surrounding
  % text (group-scheme structure on \(\IPic_{X/S}\) via tensor product);
  % arXiv:math/0504020, pp.~17, 26 (read from
  % references/kleiman-picard-src/kleiman-picard.tex, L1311-L1318 + L2057-L2069).
  % SOURCE QUOTE:
  % "\begin{dfn}\label{df:Pfs}
  %   The {\it relative Picard functor} \(\Pic_{X/S}\) is defined by
  %      \(\)\Pic_{X/S}(T):=\Pic(X_T)\big/\Pic(T).\(\)
  %   Denote its associated sheaves in the Zariski, \'etale, and fppf
  %  topologies  by
  %      \(\)\Pic_{(X/S)\zar},\ \Pic_{(X/S)\et},\
  %      \Pic_{(X/S)\fppf}.\(\)
  %   \end{dfn}"
  % SOURCE QUOTE:
  % "\begin{dfn}\label{df:Psch}
  %    If any of the four relative Picard functors of Definition~\ref{df:Pfs}
  %  is representable, then the representing scheme is called the {\it Picard
  %  scheme\/} and denoted by \(\IPic_{X/S}\).  Moreover, we say simply that the
  %  Picard scheme \(\IPic_{X/S}\) exists."
  \textit{Source: [Kleiman], ``The Picard scheme'', \S 2, Def.~2.2 +
  \S 4, Def.~4.1 (the Picard scheme).}
  Let \(\Pic_{C/k}\) be the Picard scheme of
  \cref{thm:fga_pic_representability}. Then \(\Pic_{C/k}\) is canonically an
  abelian group scheme over \(k\): the tensor product
  \([\mathcal{L}] + [\mathcal{M}] := [\mathcal{L} \otimes \mathcal{M}]\)
  and inverse \(-[\mathcal{L}] := [\mathcal{L}^{-1}]\) on
  \(\Pic^\sharp_{C/k}(T)\) (well defined because both operations
  preserve isomorphism classes and descend through the subgroup
  \(H_T = \pi_T^*\Pic(T)\) of \cref{thm:relative_pic_quotient_well_defined})
  are natural in \(T\), hence by Yoneda are represented by morphisms
  \(m : \Pic_{C/k} \times_k \Pic_{C/k} \to \Pic_{C/k}\),
  \(i : \Pic_{C/k} \to \Pic_{C/k}\), and
  \(e : \Spec k \to \Pic_{C/k}\) (the trivial line bundle class) satisfying
  the abelian-group axioms. In particular, \(\Pic_{C/k}\) is a \(k\)-group
  scheme locally of finite type.
\end{theorem}

% SOURCE QUOTE PROOF: not applicable --- Kleiman \S 2--4 do not isolate a
% separate proof block for the group-scheme structure: it is implicit in
% the fact that \(\Pic_{(X/S)\et}\) takes values in abelian groups (df:Pfs)
% and that representability transfers algebraic operations from set-valued
% functors to scheme morphisms by Yoneda. The structural restatement below
% records the standard transfer.
% NOTE (run-0008): this theorem is now PROVED in Lean, by direct Yoneda
% transport: with picSharp genuinely equal to the group-valued relative
% functor relPresheaf (def:rel_pic_sharp) with the forgetful functor
% applied, Mathlib's CommGrpObj.ofRepresentableBy applied to the
% representability witness `representable` yields the full COMMUTATIVE
% group-object structure (GrpObj + IsCommMonObj) on PicScheme C. The former
% ⟨sorry⟩-backed carrier pair (def:pic_scheme_group_object /
% def:inst_pic_scheme_group_object) is deleted. The Lean conclusion
% (CommGrpObj) is STRONGER than the abelian-group-scheme prose below asked
% for at the GrpObj level: commutativity comes with the same transport.
\begin{proof}
\leanok
  The relative Picard functor \(\Pic^\sharp_{C/k}\) of
  \cref{def:rel_pic_sharp} is abelian-group valued, with sum
  \([\mathcal{L}] + [\mathcal{M}] = [\mathcal{L} \otimes \mathcal{M}]\),
  inverse \(-[\mathcal{L}] = [\mathcal{L}^{-1}]\), and unit
  \(0 = [\mathcal{O}_{C \times_k T}]\); these operations descend through
  the quotient by \(H_T = \pi_T^*\Pic(T)\) because \(\pi_T^*\) is itself a
  group homomorphism (\cref{thm:rel_pic_addcommgroup_via_tensorobj}).

  By \cref{thm:fga_pic_representability}, \(\Pic^\sharp_{C/k}\) is
  represented by the scheme \(\Pic_{C/k}\). The Yoneda embedding is fully
  faithful and preserves limits, so a representing object of a
  group-valued presheaf is canonically a group object: the natural
  transformations
  $\Pic^\sharp_{C/k} \times \Pic^\sharp_{C/k} \to \Pic^\sharp_{C/k}$
  (multiplication), \(\Pic^\sharp_{C/k} \to \Pic^\sharp_{C/k}\) (inverse),
  and the unit map \(h_{\Spec k} \to \Pic^\sharp_{C/k}\) pick out unique
  morphisms of representing schemes \(m\), \(i\), \(e\) as displayed, and
  the abelian-group axioms hold at the level of \(T\)-points functorially
  in \(T\), hence on the representing scheme. In Lean this transport is
  exactly \texttt{CommGrpObj.ofRepresentableBy} applied to the witness of
  \cref{thm:fga_pic_representability} (transported across the
  \(\mathrm{Multiplicative}\) carrier identification); it produces the
  commutative-group-object structure directly. The local-finite-type
  property is part of \cref{thm:fga_pic_representability}.
\end{proof}

\section{The Picard scheme: the named declaration}
\label{sec:fga_pic_named_definition}

The two preceding theorems are existence statements. We record the
named declaration corresponding to the Lean target.

\begin{definition}\leanok
[The Picard scheme]
  \label{def:pic_scheme}
  \lean{AlgebraicGeometry.Scheme.PicScheme}
  % NOTE iter-283: corrected this \uses to match the actual Lean definitional
  % order and break two dependency cycles that crashed `leanblueprint web`
  % (plastexdepgraph infinite recursion). In Lean (FGAPicRepresentability.lean
  % L223) `PicScheme := (HasPicScheme.has_pic_scheme C).choose`, i.e. it is
  % extracted from the `HasPicScheme` carrier (def:has_pic_scheme) ONLY -- NOT
  % from `representable` (thm:fga_pic_representability) and NOT from the group
  % structure (thm:pic_is_group_scheme). The two removed edges were both
  % spurious back-edges:
  %   (a) thm:pic_is_group_scheme: the abelian-group-scheme structure is built
  %       ON TOP of PicScheme (def:pic_scheme_group_object -> thm:pic_is_group_scheme,
  %       both \uses def:pic_scheme), so listing it here made a 2-cycle.
  %   (b) thm:fga_pic_representability: `representable` is extracted FROM the
  %       PicSharpRepresentable carrier whose statement references PicScheme
  %       (def:pic_sharp_representable -> def:pic_scheme, thm -> def:pic_sharp_representable),
  %       so listing it here made a 3-cycle. PicScheme depends only on HasPicScheme.
  % The explanatory \cref{thm:fga_pic_representability}/\cref{thm:pic_is_group_scheme}
  % in the prose are cross-references, not dependency edges, and are retained.
  \uses{def:has_pic_scheme, def:inst_has_pic_scheme}

  % SOURCE: [Kleiman], ``The Picard scheme'', \S 4, Def.~df:Psch;
  % arXiv:math/0504020, p.~26 (read from
  % references/kleiman-picard-src/kleiman-picard.tex, L2057-L2062).
  % SOURCE QUOTE:
  % "\begin{dfn}\label{df:Psch}
  %    If any of the four relative Picard functors of Definition~\ref{df:Pfs}
  %  is representable, then the representing scheme is called the {\it Picard
  %  scheme\/} and denoted by \(\IPic_{X/S}\).  Moreover, we say simply that the
  %  Picard scheme \(\IPic_{X/S}\) exists."
  \textit{Source: [Kleiman], ``The Picard scheme'', \S 4, Def.~4.1.}
  Let \(C/k\) be a smooth proper geometrically integral curve over a field
  \(k\), with a \(k\)-rational point. The \emph{Picard scheme} \(\Pic_{C/k}\)
  is the \(k\)-scheme representing the relative Picard functor
  \(\Pic^\sharp_{C/k}\) (equivalently, by Kleiman \S 2 Thm.~2.5 under the
  rational point, its \'etale sheafification), supplied by
  \cref{thm:fga_pic_representability}, equipped with the abelian-group-scheme
  structure of \cref{thm:pic_is_group_scheme}. The scheme is separated,
  locally of finite type over \(k\), and a disjoint union of open
  quasi-projective \(k\)-subschemes.
\end{definition}

\section{Lean encoding}
\label{sec:fga_pic_lean_encoding}

The Lean target file is the new
\texttt{AlgebraicJacobian/Picard/FGAPicRepresentability.lean}, sibling to
\texttt{AlgebraicJacobian/Picard/RelativeSpec.lean}
(\cref{chap:Picard_RelativeSpec}) and
\texttt{AlgebraicJacobian/Picard/LineBundlePullback.lean}
(\cref{chap:Picard_LineBundlePullback}). The blueprint declarations of
this chapter correspond to Lean targets as follows.

\begin{itemize}
  \item \texttt{AlgebraicGeometry.Scheme.PicScheme}
    (\cref{def:pic_scheme}) is a structural constant naming the scheme
    representing \(\Pic^\sharp_{C/k}\); it is the scheme witness
    extracted from the \texttt{HasPicScheme} carrier
    (\cref{def:has_pic_scheme}) via \texttt{(HasPicScheme.has\_pic\_scheme
    C).choose}. The representability identification \texttt{representable}
    and the group-scheme structure are then layered on top of this bare
    scheme (so \cref{def:pic_scheme} depends only on \cref{def:has_pic_scheme},
    not on \cref{thm:fga_pic_representability} or
    \cref{thm:pic_is_group_scheme}).
  \item \texttt{AlgebraicGeometry.Scheme.PicScheme.picSharp}
    (\cref{def:pic_sharp_carrier}) is, after the run-0008 rewire, the REAL
    relative Picard functor: \texttt{PicSharp.relPresheaf} of
    \cref{chap:Picard_RelPicFunctor} composed with the forgetful functor
    to \(\Set\) (values in \texttt{Type (u+1)}, since the line-bundle
    carrier is a quotient of large module sheaves;
    \texttt{Functor.RepresentableBy} is universe-polymorphic, so
    representability by a scheme remains well-formed).
  \item \texttt{AlgebraicGeometry.Scheme.PicScheme.abelMap}
    (\cref{lem:line_bundle_quot_correspondence}) is the natural
    transformation \(\mathrm{Div}_{C/k} \to \Pic^\sharp_{C/k}\) sending a
    relative divisor family \(\langle \mathcal F, q\rangle\) to
    \(-[\ker q]\), the class of its associated invertible sheaf
    (run 0011: \texttt{HasAbelMap} is now data-carrying and
    \texttt{abelMap} is this substantive natural transformation, with
    defining property \texttt{abelMap\_app\_mk}; PROVED, no sorry). This
    is delivered at the level of the Type-valued carrier presheaves of
    \cref{def:div_functor_carrier,def:pic_sharp_carrier}: representability
    of the source \(\mathrm{Div}_{C/k}\) by an open subscheme of the
    Hilbert scheme (Kleiman~\S 3, Thm.~th:repDiv,
    \cref{rem:div_functor_representability}) is separate future work,
    not needed for \texttt{abelMap} itself.
  \item \texttt{AlgebraicGeometry.Scheme.PicScheme.smoothProperQuotient}
    (\cref{lem:smooth_proper_quotient}) is the Altman--Kleiman
    descent of a flat-and-proper equivalence relation \(R \subseteq Z
    \times_k Z\) on a quasi-projective \(k\)-scheme \(Z\) to a
    quasi-projective quotient \(Q\), packaged in the form that consumes
    a \(\Functor.\mathrm{RepresentableBy}\)-witness for \(Z\) and produces
    one for the coequalising sheaf, with the conclusion supplied by the
    use-site hypothesis class \cref{def:has_smooth_proper_quotient}
    (see the run-0008 note after \cref{lem:smooth_proper_quotient}).
  \item \texttt{AlgebraicGeometry.Scheme.PicScheme.representable}
    (\cref{thm:fga_pic_representability}) extracts the
    \(\Functor.\mathrm{RepresentableBy}\)-witness for
    \(\Pic^\sharp_{C/k}\) from the \texttt{PicSharpRepresentable} carrier,
    which is itself PROVED from \texttt{HasPicScheme} by
    \texttt{Exists.choose\_spec}; the four-step Kleiman assembly
    (Hilbert-polynomial stratification, \(m\)-regularity bound, Abel-map
    factorisation, smooth-proper quotient), preceded by the Kleiman \S 2
    Thm.~2.5 comparison, is the obligation of
    \cref{def:inst_has_pic_scheme}.
  \item \texttt{AlgebraicGeometry.Scheme.PicScheme.groupSchemeStructure}
    (\cref{thm:pic_is_group_scheme}) is the \texttt{CommGrpObj} instance
    on \texttt{PicScheme} obtained by Yoneda transport
    (\texttt{CommGrpObj.ofRepresentableBy}) of the abelian-group structure
    on \(\Pic^\sharp_{C/k}\) — PROVED, sorry-free, after the run-0008
    rewire.
\end{itemize}

No new core axiom is required: this chapter is an assembly chapter on
top of the QuotScheme engine, the relative-Picard \'etale sheafification,
and the line-bundle pullback machinery already in place. The single
large external input (the Altman--Kleiman theorem of effective
flat-and-proper equivalence relations on quasi-projective schemes) is
recorded as \cref{lem:smooth_proper_quotient} and is itself a standard
result; if its Lean formalisation is not yet in Mathlib, the
\texttt{smoothProperQuotient} Lean declaration may be sorried and the
gap reported to the plan agent, as it is logically downstream of the
Hilbert-scheme construction of \cref{chap:Picard_QuotScheme}.

\section{Typeclass scaffolding for the representability assembly}
\label{sec:fga_pic_typeclass_scaffolding}

The five public declarations above are each phrased so that their single
mathematical obligation is isolated in a \emph{carrier}: a \texttt{Prop}-valued
existence predicate together with a noncomputable extraction of the asserted
object. This section records, with a one-to-one Lean correspondence, each
carrier predicate, each extraction, and each existence witness. The carriers
are file-internal scaffolding for the public results; their substantive
mathematical content is the obligation analysed in
\cref{sec:fga_pic_sorry_closure_order}.

\subsection{Carriers for the sibling-chapter functors}
\label{subsec:fga_pic_functor_carriers}

\begin{definition}\leanok
[\(k\)-rational point]
  \label{def:has_rational_point}
  \lean{AlgebraicGeometry.Scheme.HasRationalPoint}
  The predicate \(\mathrm{HasRationalPoint}(C)\) asserts that the structural
  morphism \(C \to \Spec k\) admits a section
  \(\sigma : \Spec k \to C\), i.e.\ that \(C\) has a \(k\)-rational point.
  This is the standing hypothesis of the chapter (see
  \cref{sec:fga_pic_setup}): by Kleiman \S 2, Thm.~2.5, it makes the plain
  relative Picard functor \(\Pic^\sharp_{C/k}\) coincide with its \'etale
  (and fppf) sheafification, so representability can be stated against
  \cref{def:rel_pic_sharp} directly. For the pointed curves \((C, P)\) of
  the Jacobian arm the pointing \(P\) supplies this instance.
\end{definition}

\begin{proof}\leanok
  A \texttt{Prop}-valued class; defined directly in Lean.
\end{proof}

\begin{definition}\leanok
[The relative Picard functor, set-valued]
  \label{def:pic_sharp_carrier}
  \lean{AlgebraicGeometry.Scheme.PicScheme.picSharp}
  \uses{def:rel_pic_sharp}
  The set-valued shadow \(\Pic^\sharp_{C/k} : (\Sch/k)^{op} \to \Set\) of the
  group-valued relative Picard presheaf of \cref{def:rel_pic_sharp}
  (composition with the forgetful functor \(\Ab \to \Set\)).

  Run-0008 rewire: this is now a REAL definition (the former
  \texttt{Classical.choice}-of-\(\langle\)sorry\(\rangle\) carrier
  extraction, together with its typeclass \texttt{HasPicSharp} and instance
  \texttt{instHasPicSharp}, is deleted). Representability witnesses against
  it are genuine natural bijections
  \((T \to \Pic_{C/k}) \simeq \Pic(C \times_k T)/\pi_T^*\Pic(T)\).
\end{definition}

\begin{proof}\leanok
  \uses{def:rel_pic_sharp}
  Immediate from \cref{def:rel_pic_sharp}: apply the forgetful functor to
  the group-valued presheaf.
\end{proof}

\begin{definition}\leanok
[Existence carrier for the relative-divisor functor]
  \label{def:has_div_functor}
  \lean{AlgebraicGeometry.Scheme.PicScheme.HasDivFunctor}
  The predicate \(\mathrm{HasDivFunctor}(C)\) asserts that the category of presheaves
  of types on \((\Sch/k)^{op}\) admits at least one object intended to be the
  relative-effective-divisor functor \(\mathrm{Div}_{C/k}\), sending a
  \(k\)-scheme \(T\) to the set of relative effective Cartier divisors on
  \(C \times_k T\) flat over \(T\). It is the existence carrier for
  \(\mathrm{Div}_{C/k}\); since the A.2.b divisor functor
  (\cref{def:div_functor}) landed, its instance
  \cref{def:inst_has_div_functor} is proved and the class survives only as
  the pinned carrier.
\end{definition}

\begin{proof}
  Proved directly in Lean.
\end{proof}

\begin{definition}\leanok
[The relative-divisor functor of the FGA assembly]
  \label{def:div_functor_carrier}
  \lean{AlgebraicGeometry.Scheme.PicScheme.divFunctor}
  \uses{def:div_functor}
  The presheaf \(\mathrm{Div}_{C/k} : (\Sch/k)^{op} \to \Set\) is the
  relative-divisor functor \(\Div_{C/k}\) of \cref{def:div_functor},
  instantiated at the structural morphism \(C \to \Spec k\): it sends a
  \(k\)-scheme \(T\) to the set of relative effective Cartier divisors on
  \(C \times_k T\) flat over \(T\).  (Formerly an opaque chosen witness of
  \cref{def:has_div_functor}; the extraction collapsed to this direct
  definition when \cref{chap:Picard_QuotScheme} landed the divisor functor.)
\end{definition}

\begin{proof}\leanok
  \uses{def:div_functor}
  Immediate from \cref{def:div_functor}: instantiate at
  \(X = C\), \(S = \Spec k\), \(\pi\) the structural morphism.
\end{proof}

\begin{definition}\leanok
[Existence witness for the relative-divisor functor]
  \label{def:inst_has_div_functor}
  \lean{AlgebraicGeometry.Scheme.PicScheme.instHasDivFunctor}
  \uses{def:has_div_functor, def:div_functor_carrier}
  The instance recording that \cref{def:has_div_functor} holds for every
  such curve \(C/k\).  PROVED: the relative-divisor functor
  \cref{def:div_functor_carrier} witnesses the existential.  (The former
  Sorry~B of \cref{subsec:sorry_has_div_functor} is closed.)
\end{definition}

\begin{proof}\leanok
  \uses{def:div_functor_carrier}
  The functor \(\mathrm{Div}_{C/k}\) of \cref{def:div_functor_carrier} is an
  object of the presheaf category, witnessing the existential.
\end{proof}

\subsection{Carrier for the Abel map}
\label{subsec:fga_pic_abel_map_carrier}

\begin{definition}\leanok
[Data carrier for the Abel map]
  \label{def:has_abel_map}
  \lean{AlgebraicGeometry.Scheme.PicScheme.HasAbelMap,
    AlgebraicGeometry.Scheme.PicScheme.abelMap}
  \uses{def:div_functor_carrier, def:pic_sharp_carrier}
  The class \(\mathrm{HasAbelMap}(C)\) carries, as \emph{data}, a natural
  transformation of presheaves \(\mathrm{abel} \colon \mathrm{Div}_{C/k}
  \to \Pic^\sharp_{C/k}\) between the relative-divisor functor of
  \cref{def:div_functor_carrier} and the relative Picard functor of
  \cref{def:pic_sharp_carrier}; \(\mathrm{abelMap} :=
  \mathrm{HasAbelMap.abel}\) extracts it. (Repinned, run 0011, from a
  \(\mathrm{Prop}\)-valued \(\mathrm{Nonempty}\) existence gate --- whose
  \(\mathtt{Classical.choice}\) extraction carried no defining property
  --- to this data-carrying class, so that \(\mathrm{abelMap}\) now has
  the genuine defining property \Cref{lem:line_bundle_quot_correspondence}.)
\end{definition}

\begin{proof}
  Proved directly in Lean.
\end{proof}

\begin{definition}\leanok
[Substantive witness for the Abel map (Sorry~C, CLOSED)]
  \label{def:inst_has_abel_map}
  \lean{AlgebraicGeometry.Scheme.PicScheme.instHasAbelMap}
  \uses{def:has_abel_map}
  The instance recording that \cref{def:has_abel_map} holds for every such
  curve \(C/k\), populated with the \emph{substantive} witness
  \(\mathrm{abelMapWitness}(C)\) of
  \cref{lem:line_bundle_quot_correspondence} --- no longer an opaque
  choice. (The former Sorry~C of \cref{subsec:sorry_has_abel_map} is
  closed.)
\end{definition}

\begin{proof}
\leanok
  \uses{lem:line_bundle_quot_correspondence}
  \(\mathrm{abelMapWitness}(C)\) is a term of the required type
  \(\mathrm{Div}_{C/k} \to \Pic^\sharp_{C/k}\)
  (\Cref{lem:line_bundle_quot_correspondence}), witnessing the instance
  directly.
\end{proof}

\subsection{Carrier for the smooth-proper quotient}
\label{subsec:fga_pic_quotient_carrier}

\begin{definition}\leanok
[Conclusion carrier for the smooth-proper quotient lemma]
  \label{def:has_smooth_proper_quotient}
  \lean{AlgebraicGeometry.Scheme.PicScheme.HasSmoothProperQuotient}
  For a morphism \(\alpha : Z \to P\) of presheaves of types on
  \((\Sch/k)^{op}\), the predicate \(\mathrm{HasSmoothProperQuotient}(\alpha)\) asserts
  that the target \(P\) is representable. It packages the conclusion of the
  Altman--Kleiman smooth-proper quotient lemma
  \cref{lem:smooth_proper_quotient}, and is supplied at the USE SITE, where
  Kleiman's quasi-projectivity hypothesis (not expressible at the pinned
  Mathlib revision) actually holds.

  Run-0008 truth audit: the former global instance
  \texttt{instHasSmoothProperQuotient} --- a \(\langle\)sorry\(\rangle\)
  asserting this predicate for EVERY \(\alpha\), i.e.\ representability of
  every presheaf --- was mathematically false and has been deleted; see the
  note after \cref{lem:smooth_proper_quotient}.
\end{definition}

\begin{proof}\leanok
  A \texttt{Prop}-valued class; defined directly in Lean.
\end{proof}

\subsection{Carriers for the Picard scheme and its structure}
\label{subsec:fga_pic_scheme_carriers}

\begin{definition}\leanok
[Existence carrier for the Picard scheme]
  \label{def:has_pic_scheme}
  \lean{AlgebraicGeometry.Scheme.HasPicScheme}
  \uses{def:pic_sharp_carrier}
  The predicate \(\mathrm{HasPicScheme}(C)\) asserts that there exists a \(k\)-scheme
  \(X\), \emph{locally of finite type over \(k\)}, carrying a
  representability witness for the relative Picard functor
  \(\Pic^\sharp_{C/k}\) of \cref{def:pic_sharp_carrier}; in other words,
  that the Picard scheme \(\Pic_{C/k}\) exists and is locally of finite
  type. It is the existence carrier from which \cref{def:pic_scheme}
  extracts the representing object. The local-finiteness clause
  (strengthening, run 0010) is part of the same existence package as
  representability in Kleiman \S 4, Thm.~th:main --- part (1) of that
  theorem exhibits \(\Pic_{C/k}\) as a disjoint union of open
  quasi-projective \(k\)-subschemes --- so the strengthened existential is
  exactly as true as the bare one, and it lets
  \cref{def:inst_pic_scheme_lft} be proved by extraction. Its (sole,
  sorry-carrying) instance \cref{def:inst_has_pic_scheme} is conditional on
  the \(k\)-rational point \cref{def:has_rational_point}.
\end{definition}

\begin{proof}
  Proved directly in Lean.
\end{proof}

\begin{definition}\leanok
[Existence witness for the Picard scheme]
  \label{def:inst_has_pic_scheme}
  \lean{AlgebraicGeometry.Scheme.instHasPicScheme}
  \uses{def:has_pic_scheme, def:has_rational_point}
  The instance recording that \cref{def:has_pic_scheme} holds for every
  smooth proper geometrically integral curve \(C/k\) \emph{with a
  \(k\)-rational point} (\cref{def:has_rational_point}) --- run-0008: the
  conditionality is essential, since without a section the plain relative
  functor is not representable in general. This is the SINGLE genuine FGA
  sorry of the chapter after the run-0008 rewire; its discharge is the
  Step-0-through-Step-4 construction of \cref{thm:fga_pic_representability}
  (Kleiman \S 2 Thm.~2.5 + \S 4 Thm.~th:main), gated on the Riemann--Roch
  substrate. Since the run-0010 strengthening of
  \cref{def:has_pic_scheme}, this sorry also carries part (1) of Kleiman
  \S 4, Thm.~th:main: the representing scheme is locally of finite type
  over \(k\) (absorbing the former separate
  \cref{def:inst_pic_scheme_lft} sorry). See
  \cref{sec:fga_pic_sorry_closure_order}.
\end{definition}

\begin{proof}

\leanok
  Typed statement in Lean; the semantic discharge is the obligation of
  \cref{sec:fga_pic_sorry_closure_order}.
\end{proof}

\begin{definition}\leanok
[Carrier: \(\Pic_{C/k}\) represents \(\Pic^\sharp_{C/k}\)]
  \label{def:pic_sharp_representable}
  \lean{AlgebraicGeometry.Scheme.PicScheme.PicSharpRepresentable}
  \uses{def:pic_sharp_carrier, def:pic_scheme}
  The predicate \(\mathrm{PicSharpRepresentable}(C)\) asserts that the specific scheme
  \(\Pic_{C/k}\) of \cref{def:pic_scheme} is a representing object for
  \(\Pic^\sharp_{C/k}\). It is the carrier from which the representability
  witness \cref{thm:fga_pic_representability} is extracted.
\end{definition}

\begin{proof}
  Proved directly in Lean.
\end{proof}

\begin{definition}\leanok
[Existence witness for the representability identification]
  \label{def:inst_pic_sharp_representable}
  \lean{AlgebraicGeometry.Scheme.PicScheme.instPicSharpRepresentable}
  \uses{def:pic_sharp_representable, def:has_pic_scheme}
  The instance recording that \cref{def:pic_sharp_representable} holds for
  every such curve \(C/k\) (given \cref{def:has_pic_scheme}).

  Run-0008: this is now PROVED, sorry-free: \(\Pic_{C/k}\) is by definition
  the chosen witness of the \cref{def:has_pic_scheme} existential, and the
  first component of the choice specification says exactly that the chosen
  scheme represents \(\Pic^\sharp_{C/k}\) (the second component, since the
  run-0010 strengthening, is local finiteness ---
  \cref{def:inst_pic_scheme_lft}). The FGA content lives entirely upstream
  in \cref{def:inst_has_pic_scheme}.
\end{definition}

\begin{proof}\leanok
  \uses{def:has_pic_scheme}
  By the first component of \texttt{Exists.choose\_spec} applied to the
  existential of \cref{def:has_pic_scheme}.
\end{proof}

% NOTE (run-0008): the former carrier pair def:pic_scheme_group_object /
% def:inst_pic_scheme_group_object (typeclass PicSchemeGroupObject and its
% global ⟨sorry⟩ instance) is DELETED: the group-scheme structure
% thm:pic_is_group_scheme is now proved directly by Yoneda transport
% (CommGrpObj.ofRepresentableBy) from the representability witness, with no
% carrier indirection.

\begin{definition}\leanok
[Carrier: \(\Pic_{C/k}\) is locally of finite type]
  \label{def:pic_scheme_lft}
  \lean{AlgebraicGeometry.Scheme.PicScheme.PicSchemeLocallyOfFiniteType}
  \uses{def:pic_scheme, def:has_pic_scheme}
  The predicate \(\mathrm{PicSchemeLocallyOfFiniteType}(C)\) asserts that the
  structural morphism of the Picard scheme \(\Pic_{C/k}\) of
  \cref{def:pic_scheme} is locally of finite type --- part (1) of Kleiman
  \S 4, Thm.~th:main (\(\Pic_{C/k}\) is a disjoint union of open
  quasi-projective \(k\)-subschemes). It is the hypothesis the
  identity-component substrate of
  \cref{chap:Picard_IdentityComponent} consumes to specialise
  \(G^0\) to \(G = \Pic_{C/k}\). Since the run-0010 strengthening of
  \cref{def:has_pic_scheme}, its instance \cref{def:inst_pic_scheme_lft}
  is proved by extraction, and the class survives only to preserve the
  consumer signature.
\end{definition}

\begin{proof}\leanok
  A \texttt{Prop}-valued class; defined directly in Lean.
\end{proof}

\begin{definition}\leanok
[Existence witness for local finiteness]
  \label{def:inst_pic_scheme_lft}
  \lean{AlgebraicGeometry.Scheme.PicScheme.instPicSchemeLocallyOfFiniteType}
  \uses{def:pic_scheme_lft, def:has_pic_scheme}
  The instance recording that \cref{def:pic_scheme_lft} holds whenever the
  Picard scheme exists (\cref{def:has_pic_scheme}). PROVED (run 0010): the
  strengthened existential of \cref{def:has_pic_scheme} packages local
  finiteness with representability --- one existence package, as in Kleiman
  \S 4, Thm.~th:main --- so the property of the chosen witness is the
  second component of the choice specification. The hypothesis is now
  \cref{def:has_pic_scheme} rather than \cref{def:has_rational_point}: the
  rational-point conditionality lives entirely in
  \cref{def:inst_has_pic_scheme}, which supplies the
  \(\mathrm{HasPicScheme}\) instance.
\end{definition}

\begin{proof}\leanok
  \uses{def:has_pic_scheme}
  By the second component of \texttt{Exists.choose\_spec} applied to the
  existential of \cref{def:has_pic_scheme}.
\end{proof}

\section{Sorry-by-sorry closure order}
\label{sec:fga_pic_sorry_closure_order}

% Run-0008 rewrite: the iter-199 seven-sorry carrier accounting is
% superseded by the picSharp rewire. Run-0010 (fga-div front): Sorry B
% (instHasDivFunctor) is closed by the real relative-divisor functor
% (def:div_functor), and the former instPicSchemeLocallyOfFiniteType sorry
% is absorbed into Sorry A by the strengthened HasPicScheme existential.
% The historical per-sorry analyses (Sorries 1, 4, 6, 7 of the
% old numbering: instHasPicSharp, instHasSmoothProperQuotient,
% instPicSharpRepresentable, instPicSchemeGroupObject) are closed or
% deleted; see the git history of this chapter for the old text.

The Lean target file
\texttt{AlgebraicJacobian/Picard/FGAPicRepresentability.lean} carries,
after the run-0008, run-0010, and run-0011 rewires, \textbf{one} open
sorry, isolated in a single \(\langle\)sorry\(\rangle\) instance
constructor of a \texttt{Prop}-valued typeclass, so every extracted
definition and theorem body is itself axiom-clean, and consumers that
quantify over the typeclass are kernel-clean. The one:
\texttt{instHasPicScheme} (Sorry~A --- which, since the run-0010
strengthening of \cref{def:has_pic_scheme}, also carries Kleiman
th:main(1) local finiteness). The closed items of the former accounting:
the relative Picard functor carrier is REAL (\cref{def:pic_sharp_carrier}),
the relative-divisor functor carrier is REAL and its existence witness
PROVED (run 0010, \cref{def:div_functor_carrier},
\cref{def:inst_has_div_functor} --- the former Sorry~B), the substantive
Abel map is REAL and PROVED (run 0011,
\cref{lem:line_bundle_quot_correspondence}, \cref{def:inst_has_abel_map}
--- the former Sorry~C), the representability identification is PROVED
(\cref{def:inst_pic_sharp_representable}), local finiteness of the Picard
scheme is PROVED by extraction (run 0010,
\cref{def:inst_pic_scheme_lft}), the group-scheme structure is
PROVED (\cref{thm:pic_is_group_scheme}), and the false global
smooth-proper-quotient instance is DELETED
(\cref{def:has_smooth_proper_quotient}).

\subsection{Sorry A --- \texttt{instHasPicScheme} (the FGA theorem)}
\label{subsec:sorry_has_pic_scheme}

\paragraph{Location.} Instance
\texttt{AlgebraicGeometry.Scheme.instHasPicScheme}
(\cref{def:inst_has_pic_scheme}), conditional on
\texttt{[HasRationalPoint C]}.

\paragraph{Mathematical content.} The full FGA existence theorem: for a
smooth proper geometrically integral curve \(C/k\) with a \(k\)-rational
point, the relative Picard functor \(\Pic^\sharp_{C/k}\)
(\cref{def:pic_sharp_carrier}) is representable by a \(k\)-scheme. The
proof is Step~0 (Kleiman \S 2, Thm.~2.5: the section makes the plain
functor an \'etale sheaf) followed by Kleiman's four steps
(\cref{thm:fga_pic_representability}): Hilbert-polynomial stratification,
\(m\)-regularity bound, Abel-map factorisation through
\(\mathrm{Div}_{C/k}\), and the smooth-proper quotient
(\cref{lem:smooth_proper_quotient}, whose use-site
\cref{def:has_smooth_proper_quotient} instance for the Abel-map slice is
part of this obligation).

\paragraph{Blockers.} (i) the A.2.b Quot/Hilbert engine
(\cref{chap:Picard_QuotScheme}) --- the divisor functor
\(\mathrm{Div}_{C/k}\) landed in run 0010 (former Sorry~B, closed) and
the Abel map landed in run 0011 (former Sorry~C, closed;
\cref{lem:line_bundle_quot_correspondence}), but the representability
theorem Kleiman th:repDiv (\cref{rem:div_functor_representability})
remains; (ii) the Route~C Riemann--Roch substrate
(\(m\)-regularity / Serre vanishing / cohomology and base change for the
stratification, cf.\ \cref{cor:flattening_stratification_curve} and
\cref{lem:cech_flat_base_change}); (iii) Altman--Kleiman descent and
EGA~IV~8.11.5, both absent from Mathlib. This is the load-bearing sorry
of the picrep cone; it is NOT expected to close in a single session.

\paragraph{Value of the run-0008 rewire despite the open sorry.} Because
\(\Pic^\sharp_{C/k}\) is now the real functor, the extracted
\texttt{representable} witness has actual mathematical content for
consumers: quantifying over \texttt{[HasPicScheme C]}, the natural
bijection \((T \to \Pic_{C/k}) \simeq \Pic(C \times_k T)/\pi_T^*\Pic(T)\)
is available as data. The downstream tangent-space identification
(\texttt{Pic0.tangentSpaceIso}, chapter
\texttt{Picard\_Pic0AbelianVariety}) attaches to exactly this bijection
with \(T = \Spec k[\epsilon]\); it was impossible against the former
opaque carrier.

\subsection{Sorry B --- \texttt{instHasDivFunctor} (CLOSED, run 0010)}
\label{subsec:sorry_has_div_functor}

\paragraph{Location.} Instance
\texttt{AlgebraicGeometry.Scheme.PicScheme.instHasDivFunctor}
(\cref{def:inst_has_div_functor}).

\paragraph{Status.} CLOSED (run 0010): \cref{chap:Picard_QuotScheme}
landed the relative-divisor functor \(\Div_{X/S}\)
(\cref{def:div_functor} --- the subfunctor of the Hilbert functor
\(\Hilb_{C/k} = \Quot_{\mathcal{O}_C/C/k}\) of relative effective Cartier
divisors, in the invertible-kernel quotient encoding of
\cref{def:div_family}), the carrier \(\mathrm{Div}_{C/k}\)
(\cref{def:div_functor_carrier}) is its instantiation at
\(C \to \Spec k\), and the instance is proved with that witness. The
base-change well-definedness of the divisor condition remains a typed leaf
(\cref{lem:relative_divisor_base_change}); representability of
\(\Div_{C/k}\) by an open subscheme of the Hilbert scheme (Kleiman \S 3,
Thm.~th:repDiv) is future work, not consumed by this chapter.

\subsection{Sorry C --- \texttt{instHasAbelMap} (CLOSED, run 0011)}
\label{subsec:sorry_has_abel_map}

\paragraph{Location.} Instance
\texttt{AlgebraicGeometry.Scheme.PicScheme.instHasAbelMap}
(\cref{def:inst_has_abel_map}).

\paragraph{Status.} CLOSED (run 0011): \texttt{HasAbelMap} is repinned
from a \texttt{Prop}-valued \texttt{Nonempty} existence gate (whose
\texttt{Classical.choice} extraction carried no defining property) to a
\emph{data-carrying} class with field \texttt{abel}
(\cref{def:has_abel_map}). The field is populated by the substantive
witness \(\mathrm{abelMapWitness}(C) = \langle \mathcal F, q\rangle
\mapsto -[\ker q]\) of \cref{lem:line_bundle_quot_correspondence}: the
composite of the kernel-class natural transformation
\(\mathrm{abelKernelNatTrans}(C)\) (naturality via the kernel--flat-base-change
comparison isomorphism) with the group-presheaf negation
\(\mathrm{picNeg}(C)\). Consequently \(\mathrm{abelMap} :=
\mathrm{HasAbelMap.abel}\) has the genuine defining property
\texttt{abelMap\_app\_mk} (proved by \texttt{rfl}), rather than being an
opaque \texttt{Classical.choice} of a \texttt{Nonempty}. As in Sorry~B,
representability of \(\mathrm{Div}_{C/k}\) by an open subscheme of the
Hilbert scheme (Kleiman th:repDiv) is \emph{not} needed or claimed here
(\cref{rem:div_functor_representability}); the Abel map is delivered at
the level of the Type-valued carrier presheaves.

\subsection{Closure order}
\label{subsec:fga_pic_closure_order_summary}

Dispatch order: \(B \to C \to A\), with \(B\) CLOSED in run 0010 (the
A.2.b divisor functor landed) and \(C\) CLOSED in run 0011 (the
substantive Abel map \(\langle \mathcal F, q\rangle \mapsto
-[\ker q]\)). Sorry~A additionally needs the Route~C Riemann--Roch
substrate and the Altman--Kleiman/EGA descent inputs, and consumes both
\(B\) and \(C\) (and, since run 0010, also discharges Kleiman th:main(1)
local finiteness). Downstream Route~A consumers proceed under typeclass
quantification (\texttt{[HasPicScheme C]}) and stay kernel-clean; the
\texttt{[HasRationalPoint C]} hypothesis is supplied by the pointing
\(P : \mathbf 1 \to C\) already present in the Jacobian statements.

\section{Out of scope}
\label{sec:fga_pic_out_of_scope}

This chapter packages \emph{only} the representable Picard scheme
\(\Pic_{C/k}\) as an assembly of sibling-chapter inputs, plus the
group-scheme refinement. The following downstream constructions live in
separate sibling chapters and are explicitly out of scope here:
\begin{itemize}
  \item \textbf{Identity component \(\Pic^0_{C/k}\)} (A.3) --- the open and
    closed subgroup scheme of \(\Pic_{C/k}\) containing the unit and
    consisting of those classes lying in the connected component of
    \(\Pic_{C/k}\). Its dimension, smoothness, and abelian-variety
    structure are the content of the A.3 chapter
    \texttt{Picard\_Pic0AbelianVariety.tex} (planned), gated on this
    chapter.
  \item \textbf{Quot scheme representability} --- consumed by
    \cref{thm:fga_pic_representability} via \texttt{thm:quot\_representable}
    (\cref{chap:Picard_QuotScheme}), not re-proved here.
  \item \textbf{\'Etale sheafification of \(\Pic^\sharp\)} ---
    \cref{chap:Picard_RelPicFunctor} (A.1.c) constructs the relative
    Picard presheaf; a topology-parametric sheafification
    (\texttt{PicSharp.etSheaf}) exists there but, after the run-0008
    route-(b) rewire, this chapter does NOT consume it: the standing
    rational-point hypothesis makes the plain relative functor already an
    \'etale sheaf (Kleiman \S 2, Thm.~2.5).
  \item \textbf{Universal Poincar\'e sheaf} (Kleiman~\S 4, Ex.~4.3)
    --- exists when \(\Pic_{C/k}\) represents \(\Pic_{C/k}\) itself and
    \(C/k\) has a section. Under the chapter's standing rational-point
    hypothesis both conditions in fact hold (Thm.~2.5 identifies the
    plain functor with the represented sheaf), so the Poincar\'e sheaf
    becomes available downstream --- but it is not constructed or
    consumed in this chapter.
  \item \textbf{Mumford's example} (Kleiman~\S 4, Eg.~4.14), counter-examples
    to representability for non-integral geometric fibers --- not
    relevant to a smooth proper geometrically integral curve.
\end{itemize}

\section{Internal-consistency check}
\label{sec:fga_pic_consistency_check}

The five declaration blocks form a connected \(\uses\) chain rooted in the
sibling chapters \cref{chap:Picard_LineBundlePullback},
\cref{chap:Picard_QuotScheme}, and \cref{chap:Picard_RelPicFunctor}:
\begin{itemize}
  \item \cref{lem:line_bundle_quot_correspondence} uses
    \cref{def:line_bundle_on_product} (A.1.b) and
    \cref{thm:pullback_natural} (A.1.b).
  \item \cref{lem:smooth_proper_quotient} is a structural Archon-internal
    statement; no external \texttt{\textbackslash uses} pointers.
  \item \cref{thm:fga_pic_representability} uses
    \cref{lem:line_bundle_quot_correspondence},
    \texttt{thm:quot\_representable} (A.2.b),
    \texttt{def:rel\_pic\_sharp} (A.1.c),
    \cref{def:has_rational_point},
    \cref{thm:relative_pic_quotient_well_defined} (A.1.b), and
    \cref{lem:smooth_proper_quotient}.
  \item \cref{thm:pic_is_group_scheme} uses
    \cref{def:pic_scheme}, \cref{thm:fga_pic_representability},
    \texttt{def:rel\_pic\_sharp} (A.1.c), and
    \cref{thm:relative_pic_quotient_well_defined} (A.1.b).
  \item \cref{def:pic_scheme} uses
    \cref{def:has_pic_scheme} and \cref{def:inst_has_pic_scheme} only
    (see the iter-283 cycle-breaking note at the definition).
\end{itemize}
The label \texttt{thm:quot\_representable} is a forward reference to the
sibling chapter \texttt{Picard\_QuotScheme.tex};
\texttt{def:rel\_pic\_sharp} is the (run-0008 repinned) relative Picard
functor of \texttt{Picard\_RelPicFunctor.tex}. If either changes label,
the \texttt{\textbackslash uses} pointers should be regenerated across
the chapters in lockstep.
